\documentclass[10pt]{article}
\usepackage[a4paper,margin=1.7cm]{geometry}
\usepackage{amsmath,amssymb,booktabs,array,enumitem,xcolor}
\usepackage[hidelinks]{hyperref}
\setlist{nosep,leftmargin=1.2em}
\newcommand{\solved}{\textcolor{green!45!black}{\textbf{solved}}}
\newcommand{\partl}{\textcolor{orange!80!black}{\textbf{partial}}}
\newcommand{\prog}{\textcolor{gray}{\textbf{in progress}}}
\begin{document}
\begin{center}
{\Large\bfseries Climbing to the Frontier --- RESULTS}\\[2pt]
BAINSA Hackathon, Track 3 (Anthropic), 26 September 2026\\
Team: Lorenzo Pazienza (team lead), Seby, Teo, Niccol\`o
\end{center}

\noindent\textbf{In one line.} We hand in a write-up for 22 of the 24 cells (P1~C6 and P3~C6 were still being written at the time of this document; see the portal). 16 cells are complete proofs (P2~C5 counts here: both of its bounds are proved), 4 are partial and say exactly which part is proved, and 2 (P4~C4/C5) are computational certificates. Every cell file states \emph{solved} or \emph{partial} on its first page. Anything proved only with a computer is labelled \emph{computer-assisted} and ships with its code and log.

\medskip
\begin{tabular}{@{}l l p{12.2cm}@{}}
\toprule
Cell & Status & What we claim\\
\midrule
\multicolumn{3}{@{}l}{\textbf{P1 --- Angles between lines (Fejes T\'oth)}}\\
C1 & \solved & Planar case, self-contained proof (the result itself is known [FVZ16, BM19]).\\
C2 & \solved & The orthogonality lemma.\\
C3 & \solved & $d+1$ lines in $\mathbb R^d$, via a saturation step and a matching argument.\\
C4 & \solved & 5 lines in $\mathbb R^3$, 6 in $\mathbb R^4$. The 6-line case is also given a second proof via the C5 cycle lemma.\\
C5 & \solved & $N=d+2$ lines for \emph{every} $d$. New ingredient: a cycle lemma (tridiagonal/cyclic nullity plus a monotone potential), checked numerically on $4\cdot10^4$ cycles.\\
C6 & \prog & Partial: $N=6$ in $\mathbb R^3$ by subset averaging; structural reduction for $N=d+3$.\\
\midrule
\multicolumn{3}{@{}l}{\textbf{P2 --- Uphill paths on the hypercube}}\\
C1--C4 & \solved & $U(Q_3)=14$, $U(Q_4)=34$, $U(Q_5)=88$, $U(Q_6)=204$, $U(Q_7)=464$, $U(Q_8)=1040$, all \emph{with proofs of optimality}. Key new tool: an exact formula $P(f)=|V|+\sum_{u\in R}(\deg u-1)+\sum_{u\in R}(N(u)-2)\,\mathrm{out}(u)$, which gives $U(G)\ge |V|+(d-1)\nabla(G)$ for $d$-regular $G$ ($\nabla$ = decycling number). Code labellings give the matching upper bound. $\nabla(Q_5)\ge14$ and $\nabla(Q_6)\ge28$ are SAT facts, each checked with two solvers and by a second, independently written program.\\
C5 & \solved\ (bounds) & $\mathbf{2328\le U(Q_9)\le 2400}$. The lower bound is new: it beats the counting bound 2312, using a split of $Q_9$ into eight $Q_6$ copies, a parity lemma (SAT) and a Delsarte LP bound. The upper bound comes from a 20-word even code.\\
C6 & \partl & Reduction: beating 2400 requires a decycling set of $Q_9$ of size $\le235$. Conjecture: $U(Q_9)=2400$. A direct SAT search for such a set ran for about 3 hours without an answer.\\
\midrule
\multicolumn{3}{@{}l}{\textbf{P3 --- Bulgarian solitaire}}\\
C1, C2 & \solved & Cyclic partitions and cycles; $D_B(T_k)$ (reproduces Griggs--Ho, with complete proofs).\\
C3 & \partl & The general bound and $D_B(T_k-1)$ are proved. The extremal set is given only by an exact dynamical criterion.\\
C4 & \solved & $D_B(T_{k-1}+1)=(k-1)(k-3)$ for all $k\ge5$, with no computation. This is the $r=1$ case of Griggs--Ho Conjecture 4.7, which we did not find proved in the literature we checked.\\
C5 & \partl & $D_B(T_{k-1}+2)=(k-1)(k-4)$ for all $k\ge7$ by one uniform argument (Conj.~4.7, $r=2$), and values $2,3,5,8,12$ for $k=2..6$. One explicit extremal family is given; the full extremal set only via a criterion.\\
C6 & \prog & Partial: exact values against Conjecture 4.7 for small $n$; the proved families from C2--C5.\\
\midrule
\multicolumn{3}{@{}l}{\textbf{P4 --- Disjoint congruence classes (Sun's conjecture)}}\\
C1--C3 & \solved & $k=3$, $k=4$, $k\le 8$ (known from O'Bryant; independent proofs).\\
C4 & \solved & All $k\le12$. For each $9\le k\le16$ we report exactly what is left undecided, which is nothing: every survivor is decided UNSAT.\\
C5 & \solved\ (cert.) & Certified boundary $K^\ast=15$, meeting requirements (i)--(v): the vacuity test finds genuine systems; every survivor has an individual decision; exact survivor lists with SHA-256; the extra pruning rule is proved and switching it off leaves the lists unchanged; two independent implementations (SAT and DFS) give elementwise-identical lists for all $k\le16$. $k=16$ is decided by the SAT implementation alone; the DFS cross-check is partial, so we do not claim 16. $k=24$ and $k=30$ are decided by hand (the prime $k-1$ argument).\\
C6 & \partl & Part (a), ``decide a size $k\ge25$'': $k=32$, $38$ and $42$ are decided, computer-assisted (O'Bryant's method with new ingredients, finishing with an exhaustive 4\,s check over a finite set of prime distributions, rerun reproduced). $k=44$ is computer-assisted but the certificate is incomplete (runtime estimate $\approx240$ CPU-hours). $k=48$ is open.\\
\bottomrule
\end{tabular}

\medskip
\noindent\textbf{Cells 5 and 6.} All four C5 cells carry a real result: P1~C5 for all $d$; a new lower bound in P2~C5; the $r=2$ case of the Griggs--Ho conjecture in P3~C5; and a certified boundary with two independent implementations in P4~C5. Of the open C6 cells, we hand in partial progress on P2 and P4, and on P1 and P3 as they are completed.

\medskip
\noindent\textbf{What we do \emph{not} claim.} We do not claim the full Fejes T\'oth conjecture, the exact value of $U(Q_9)$, the full function $D_B(n)$, $k=16$ in the certified range, or $k\ge44$ for Sun's conjecture. ``New'' means new relative to the sources we cite (Griggs--Ho 1998, O'Bryant 2007, Bilyk--Matzke 2019); we did not survey the literature exhaustively.

\medskip
\noindent\textbf{Where to find it.} One folder per cell (\texttt{P\emph{i}\_C\emph{j}/}) with the PDF, the \LaTeX\ source, code and logs; \texttt{INDEX.md} lists the file handed in for every cell.
\end{document}
